\documentclass[11pt,a4paper]{article}
\usepackage[margin=1.9cm]{geometry}
\usepackage{amsmath,amssymb,booktabs,enumitem}
\usepackage[expansion=false]{microtype}
\usepackage[hidelinks]{hyperref}
\usepackage{xcolor}
\linespread{0.98}
\newcommand{\code}[1]{\texttt{\small #1}}
\newcommand{\flag}[1]{\textcolor{red!70!black}{#1}}
\begin{document}
\begin{center}
{\large\bfseries Novelty of the RadCluster cluster-dynamics method}\\[2pt]
{\footnotesize Verified against \code{Ghoniem-Org/RadCluster}, \code{RadCluster\_2\_1};
benchmarked against Kohnert \& Wirth, \emph{Modelling Simul.\ Mater.\ Sci.\ Eng.} \textbf{25} (2017) 015008 --- \today}
\end{center}
\vspace{2pt}

\section{Production bias and simultaneous loop--cavity evolution}

\paragraph{The source term.}
Cascade production (\code{py\_utils/defect\_production.py}) is a power-law
partition, for $x=i,v$,
\[
G^x_1=\eta\dot d\,(1-f^x_{\rm cl}),\qquad
G^x_n=\eta\dot d\,A_x n^{-s_x}\ (n\ge2),\qquad
\sum_n n\,G^x_n=\eta\dot d ,
\]
the monomer channel absorbing the exact remainder so that atom conservation holds
even under grid truncation. \emph{Equal total SIA and vacancy inventories are
created; they are assigned to different mobility classes.} That asymmetry is the
production bias. For the fission spectrum:

\begin{center}\footnotesize\setlength{\tabcolsep}{6pt}
\begin{tabular}{lccc}
\toprule
Measure & SIA & vacancy & ratio \\
\midrule
clustered atom fraction & 0.580 & 0.150 & 3.87 (i/v) \\
cluster-number fraction & 0.1086 & 0.0485 & \textbf{2.24} (i/v) \\
monomer atom fraction & 0.420 & 0.850 & \textbf{2.02} (v/i) \\
mean produced cluster size & 5.34 & 3.09 & --- \\
\bottomrule
\end{tabular}
\end{center}

\noindent The ``factor of two'' is correct for produced \emph{cluster number}, and
equivalently for the inverse \emph{monomer} ratio. The clustered atom partition
differs by 3.9. Quote whichever is load-bearing, but not ``2'' for the atom
fraction.

\paragraph{What the model does with it.}
Cascade-born SIA clusters seed loops; the larger vacancy-monomer source feeds
cavities. RadCluster evolves both ladders with SIA--SIA and V--V coalescence,
residual-size reactions $I_n+V_m\rightarrow I_{n-m},\,V_{m-n}$ or $\varnothing$,
cavity capture and emission, helium processes, and sink accounting. The
size-resolved \emph{coupled} evolution of both families is the substantive
advance over a two-monomer swelling model.

The mixed-dimensional cavity kernel (\code{py\_utils/reaction\_rates.py}) is
\[
\mathcal{K}^{\rm cav}_{n,m}\;\propto\;
\frac{m^{1/3}\,\omega^{\rm 1D}_n}{1+B_{\rm rot}\hat L^2 m^{-1/3}},
\qquad B_{\rm rot}=\tfrac{4}{\pi}\bigl(\tfrac{8\pi}{3}\bigr)^{1/3}\!\approx2.627 .
\]
\flag{Do not quote the suppression factor without $m$.} The often-cited
$1+B_{\rm rot}\hat L^2$ is the $m\!=\!1$ reference; larger cavities are
progressively \emph{less} suppressed through $m^{-1/3}$. The defensible statement
is that \emph{dimensionality changes the relative rates of cavity interception,
coalescence and fixed-sink capture} --- not that 1D clusters cannot find
cavities. Likewise, coalescence reduces mobile cluster number and transfers
inventory into large or sessile loops, but it \emph{conserves} SIA mass; it is a
redistribution between mobility classes, not a sink.

\paragraph{The exact balance.}
With balanced production and initial inventories, the code enforces
(\code{py\_utils/post\_process.py}, $\delta_{\rm FP}<10^{-6}$)
\begin{equation}
S_V(t)-S_I(t)\;=\;J^{\rm sink}_I(t)-J^{\rm sink}_V(t),
\label{eq:balance}
\end{equation}
up to numerical residuals and boundary losses. Internal coalescence and mutual
annihilation cancel identically: both destroy or relocate equal numbers of the
two species.

\paragraph{\flag{Correction: equal capture efficiencies do not zero the right-hand side.}}
An earlier draft of these notes proposed setting $Z_i=Z_v$ to force
$\Delta J\!\to\!0$ and collapse \eqref{eq:balance} to $S_V=S_I$. \textbf{That is
wrong.} Equal $Z$ removes \emph{preferential} elastic capture, but the sink
fluxes remain unequal because the mobile concentrations and diffusivities of the
two species differ --- and they differ \emph{precisely because of} the production
partition. Equal $Z$ therefore does not force $S_V=S_I$; it removes one cause of
the imbalance, not the imbalance.

\paragraph{The decisive design.}
Isolation requires a $2\times2$ matrix, not a single run:

\begin{center}\footnotesize
\begin{tabular}{l|cc}
& $f^i_{\rm cl}=f^v_{\rm cl}$ & $f^i_{\rm cl}>f^v_{\rm cl}$ \\
\midrule
$Z_i=Z_v$    & reference        & \textbf{production bias only} \\
$Z_i\neq Z_v$ & sink bias only   & combined (shipped EUROFER deck) \\
\end{tabular}
\end{center}

\noindent Equalise the dislocation, grain-boundary and precipitate capture
efficiencies together ($Z^d$, $Z^{gb}$, $Z^p$ --- the shipped deck has
$Z_i/Z_v=1.35$, $Z^{gb}_{i,v}=1.40/0.93$, $Z^{p}_{i,v}=1.45/0.93$ at
$\rho_d=5\times10^{14}\,\mathrm{m^{-2}}$) while holding sink densities,
transport, binding, helium conditions and dose rate fixed. The claim is supported
if, in the production-only cell, \emph{both} loop and cavity number densities and
mean sizes grow relative to the unbiased-source reference. Report the source,
coalescence, annihilation, cavity-capture and external-sink fluxes against
\eqref{eq:balance} for each cell.

\paragraph{What is, and is not, new.}
Golubov--Singh--Trinkaus rate theory already \emph{permits} production-biased
swelling; that is not the novelty. The novelty is the simultaneous
\emph{size-resolved} CD evolution of both cluster families under mixed-dimensional
transport and nonlocal (residual-size, coalescence) reactions, with the two bias
mechanisms separated by controlled construction rather than by asymptotic argument.

\section{Bin-moment reduction versus existing grouping techniques}

RadCluster stores $\mu^{(p)}_k=\sum_{n\in\mathcal{B}_k}n^p c_n$ and implements
\[
\{\mu^{(p)}_k\}\ \xrightarrow{\ \textbf{R}\ }\ c^{\rm rec}_n
\ \xrightarrow{\ \textbf{E}\ }\ \dot c^{\rm full}_n
\ \xrightarrow{\ \textbf{P}\ }\ \dot\mu^{(p)}_k=\sum_{n\in\mathcal{B}_k}n^p\dot c^{\rm full}_n ,
\]
reconstructing, evaluating the \emph{per-size} master equation, and projecting --- at
every RHS call (\code{core/reductions/bin\_moment.py}).

\paragraph{Kiritani (1973).} One average concentration per group with
group-averaged rate constants. Inconsistent: densities and their increments are
assumed equal within a group, breaking $\mathrm{d}N/\mathrm{d}t$ and
$\mathrm{d}S/\mathrm{d}t$.

\paragraph{Golubov \emph{et al.} (2001).} Linear intra-group profile
$f(x)=L^0_i+L^1_i(x-\langle x\rangle_i)$ with both moments evolved; number and
content are conserved by construction. Not intrinsically nonconservative --- the
limitations are that it was formulated for \emph{monomer} flux only, and that the
intra-group mean flux $\langle J\rangle^*_i$ is undefined at $\Delta x_i=1$.

\paragraph{Kohnert \& Wirth (2017) --- the real benchmark.} They remove Golubov's
monomer restriction, generalising the moment scheme to arbitrary cluster--cluster
reactions, and add four interpolation schemes (linear, secant, Lagrange, cubic
spline). Both families work by \emph{projecting the rate constants onto basis
functions} to build grouped coefficients $a^{\pm},b^{\pm}$ (their Eqs.~7, 19, 20
and Appendix~A), which are ``independent of the actual concentration values so
they need be computed only once, before time integration begins.''

That precomputation is what buys their speedup, and it carries an explicit price,
stated twice in their paper: ``\emph{Enforcing constant cluster properties across
each group} allows the rate constants to be extracted from the summations'' (p.~4),
and ``due to the assumption that rate constants are unchanging across a group,
clusters whose properties differ radically from neighbors in phase space should
not be grouped over\ldots we enforce a minimum size to the grouped region\ldots
typically around clusters with 20 members'' (p.~9).

\paragraph{The distinction.} R--E--P forms \emph{no grouped operator at all}, so
the constant-property assumption is never made: step \textbf{E} evaluates
$\mathcal{K}(n,n')$ at the true size. This matters for the physics of \S1
specifically, because the fastest-varying kernel in the model is the one carrying
the result: $\mathcal{K}^{\rm cav}_{n,m}$ depends on $n$ through
$\omega^{\rm 1D}_n\propto n^{-s}$ and on $m$ through
$m^{1/3}/(1+B_{\rm rot}\hat L^2m^{-1/3})$, both strongly non-linear across a
logarithmic bin. A scheme that holds cluster properties constant over the group
flattens exactly the quantity that sets the relative rates of cavity interception
and sink capture. \textbf{Reaction fidelity, not work reduction, is the advantage.}

A second practical consequence: because no reaction-specific grouped operator is
derived, a new channel costs no new coefficient algebra. Kohnert \& Wirth report
that setup ``becomes prohibitive, particularly when applying the moment scheme
for problems with dense reaction networks,'' requiring computer-algebra-derived
analytic forms per reaction class. RadCluster's loop$\rightarrow$network loss and
the planned RIS/precipitation edges reuse the existing reduction and conservation
machinery unchanged.

\begin{center}\footnotesize\setlength{\tabcolsep}{4pt}
\begin{tabular}{@{}p{1.9cm}p{2.7cm}p{3.0cm}p{3.5cm}p{4.3cm}@{}}
\toprule
& \textbf{Kiritani (1973)} & \textbf{Golubov (2001)} & \textbf{Kohnert--Wirth (2017)} & \textbf{RadCluster (R--E--P)} \\
\midrule
State & uniform group density & $L^0,L^1$; linear profile & interpolation nodes \emph{or} $L^0,L^1$ & constant / linear / log-normal, $P=1,2,3$ \\
Operator & grouped equation & derived group equations & \textbf{precomputed} grouped network & reconstruct grid, per-size RHS, project \\
Kernels & group approximations & grouped coefficients & grouped; \textbf{constant across group} & \textbf{original size-dependent kernels} \\
Reactions & monomer & monomer only & arbitrary cluster--cluster & arbitrary; same machinery \\
New channel & --- & re-derive & new analytic coefficients & \textbf{none required} \\
Conservation & may lose consistency & $N,S$ by construction & moments by construction; interp.\ leaks $O(\Delta x^2)$ & projected stoichiometry, audited \\
Grouped range & --- & --- & $\gtrsim20$ members & no lower restriction \\
Work & group-level & group-level & \textbf{group-level (cost $\downarrow$ with width)} & reduced state, full-grid sweep $I{+}V$ \\
\bottomrule
\end{tabular}
\end{center}

\paragraph{\flag{Qualifications that must accompany the claim.}}
Conservation is \emph{not} unconditional. Linear reconstruction clips negative
values (\code{np.maximum} in \code{distribution\_from\_moments\_hat}); the discrete
log-normal closure normalises $\mu^{(0)}$ exactly but recovers higher moments only
approximately; the upper grid boundary can lose products; solver tolerances and
floors enter the residual. The correct claim is that RadCluster is
\emph{stoichiometrically conservative at the projected-reaction level and
continuously audited} via $\delta_{\rm FP}$, $\delta_{\rm He}$. Kohnert \& Wirth
report the analogous pathologies for their schemes --- SDF oscillation, negative
concentrations, phase-space-boundary failure --- so this is a shared hazard, not a
differentiator. Finally, the $O((r{-}1)^2)$, $O((r{-}1)^3)$, $O((r{-}1)^4)$
truncation orders should be removed unless derived for the implemented
\emph{discrete} closures, and R--E--P forgoes the asymptotic work reduction that is
Kohnert \& Wirth's headline result.

\section*{Contribution statement}

\emph{\small
RadCluster is a production-spectrum-resolved cluster-dynamics framework that
couples asymmetric cascade partitioning, mixed 1D/3D SIA transport, SIA--SIA and
vacancy--vacancy coalescence, residual-size SIA--vacancy reactions, cavity capture
and emission, helium coupling, and explicit sink accounting, so that loop and
cavity populations evolve as size-resolved distributions in a single coupled
system rather than as separately nucleated averages. On the physics side,
controlled calculations that independently equalise cascade clustering fractions
and sink capture efficiencies isolate simultaneous loop and cavity growth caused
by production bias rather than by preferential sink capture --- a separation that
reduced-moment rate theory can argue but not demonstrate. On the numerical side,
its reconstruct--evaluate--project bin-moment reduction retains the original
size-dependent reaction kernels instead of projecting them onto grouped
coefficients, so no constant-property assumption is imposed within a bin and no
reaction-specific inter-group operator need be derived for new physics; the
reduction conserves projected stoichiometric balances under continuous audit of
reconstruction, boundary and solver error, trading the asymptotic work reduction
of a strict grouped operator for reaction fidelity in exactly the size-dependent,
mixed-dimensional kernels that govern the swelling result.}

\end{document}
